\subsection{Stable scaled power-series arithmetic}
\label{sec:scaled-series}

\begin{lemma}[Coefficient computation with rounded arithmetic]
\label{lem:scaled-series-stability}
Consider a zero-bias tanh network with $D$ layers, width $n$, and
absolute row sums at most one. Let each input be an affine polynomial
in a complex variable $u$. Suppose that on $|u|\le2$ all exact neuron
values and preactivations have modulus at most two, that
$|e^{2s}|\le e^4$, and that $|(1+e^{2s})^{-1}|\le1$ for every
preactivation $s$. For each neuron, all Taylor coefficients through
degree $K$ can be enclosed to width $2^{-S}$ using
\[
 O\bigl(D[n^2K+nK^2+nM]M^2\bigr)
\]
bit operations, where it suffices to take
\[
 M=S+64D+\lceil\log_2(256nD(K+2)^6)\rceil+b+64.
\]
The coefficients are stored in the scaled variable $u$; there is no
inverse power of the original disk radius in this bound.
\end{lemma}
\paragraph{Proof idea.}
Work with coefficients in a scaled variable and their sum norm. Residual
identities for the exponential and reciprocal recurrences control local
rounding, while the row norm bounds propagation between layers.

\begin{proof}
For truncated series write $\|A\|=\sum_{j=0}^K|[u^j]A|$.
This is a submultiplicative norm on the quotient algebra modulo
$u^{K+1}$. Cauchy's estimate on the radius-two disk bounds the
exact coefficient norms of neurons, preactivations, exponentials,
and inverses respectively by $4,4,2e^4,2$.

Round each coefficient operation outward to $M$ fractional bits.
For $E=e^A$, compute the constant coefficient by a certified
bounded-argument exponential and use
\[
 e_0=e^{a_0},\qquad
 k e_k=\sum_{j=1}^k j a_j e_{k-j}\quad(1\le k\le K).
\]
Each coefficient equation has at most $K$ rounded products with
integer factor at most $K$. Its residual in $E'-A'E$ has norm at
most $4(K+2)^3 2^{-M}$. Solving this residual equation by
multiplication by $e^{-A}$, coefficientwise integration, and
multiplication by $e^A$ gives total error at most
$2^{32}(K+2)^4 2^{-M}$ whenever $\|A\|\le9$.
Indeed both exponential norms are at most $e^9$, and integration
has norm at most one; the initial constant error fits this bound.
The inverse is computed from its triangular product recurrence.
Its product residual has the same polynomial bound. Multiplication
by the exact inverse bounds the error; a perturbation of norm below
$1/4$ keeps the perturbed inverse norm at most four.

For propagation between rows use
\[
 \|e^{A+E}-e^A\|\le e^{\|A\|+\|E\|}\|E\|,\qquad
 (U+E)^{-1}-U^{-1}=-(U+E)^{-1}EU^{-1}.
\]
Since $\|2s\|\le8$, these inequalities and
$\tanh s=1-2/(1+e^{2s})$ bound error amplification through a
nonlinear layer by $2^{20}$. A linear row does not amplify the
incoming norm error, since its absolute row sum is at most one.
Local rounded linear sums and nonlinear operations contribute at
most $2^{32}n(K+2)^4 2^{-M}$. Induction therefore bounds the error
at any layer by
\[
 D2^{20D+32}n(K+2)^4 2^{-M}.
\]
The displayed choice of $M$ makes this smaller than $2^{-S-4}$,
and also verifies inductively each $1/4$ perturbation margin used
above. Adding rational outward error bounds gives intervals of
width at most $2^{-S}$. The initial affine coefficients and dyadic
weights are represented exactly, or enclosed with the same error
budget if division is needed to describe the affine input.

There are $O(n^2DK)$ linear coefficient operations and
$O(nDK^2)$ nonlinear coefficient operations. Under the norm bounds,
integer parts have $O(M)$ bits. Schoolbook arithmetic costs
$O(M^2)$ per operation. Each constant exponential has
$O(M^3)$ bit cost by a rounded Taylor expansion and bounded range
reduction. These counts give the claimed bound.
\end{proof}
